\subsubsection{Analisis Komponen Infrastruktur Audit Trail}

Setelah kebutuhan \textit{audit event} ditentukan, langkah selanjutnya adalah mengidentifikasi komponen infrastruktur yang diperlukan agar proses pencatatan, penyimpanan, pemantauan, dan analisis \textit{audit trail} dapat berjalan secara menyeluruh. Analisis ini dilakukan dengan mengacu pada model konseptual sistem \textit{audit trail} yang dijelaskan pada Annex B ISO 27789. Berbeda dengan NIST SP 800-53 yang mendefinisikan kontrol keamanan yang harus dipenuhi oleh suatu sistem \textit{audit trail}, ISO 27789 menyediakan model layanan (\textit{service-oriented architecture}) yang menggambarkan komponen-komponen utama beserta interaksinya dalam membangun sebuah sistem \textit{audit trail}.

Model tersebut menunjukkan bahwa sistem \textit{audit trail} tidak hanya terdiri atas media penyimpanan log, tetapi merupakan sekumpulan layanan yang bekerja sama untuk menerima \textit{audit event}, membentuk \textit{audit record}, menyimpan data, melakukan pemantauan terhadap aktivitas yang terjadi, hingga menyediakan fasilitas pencarian dan pelaporan bagi auditor. Berdasarkan model tersebut, komponen-komponen yang menjadi kebutuhan dasar infrastruktur \textit{audit trail} dapat diidentifikasi sebagaimana ditunjukkan pada Tabel \ref{tab:komponen_audit_trail_iso}.

\begin{xltabular}{\textwidth}{|p{4cm}|X|}
\caption{Identifikasi Komponen Infrastruktur Audit Trail berdasarkan ISO 27789 Annex B}
\label{tab:komponen_audit_trail_iso}\\

\hline
\textbf{Komponen} & \textbf{Fungsi} \\
\hline
\endfirsthead

\hline
\textbf{Komponen} & \textbf{Fungsi} \\
\hline
\endhead

\hline
\endfoot

Audit Logger  &
Menerima \textit{audit event} dari aplikasi atau sistem, melakukan validasi awal, kemudian meneruskan \textit{event} untuk diproses dan disimpan sebagai \textit{audit record}. \\
\hline

Audit Record Generator  &
Membentuk \textit{audit record} berdasarkan jenis \textit{audit event} yang diterima sehingga seluruh \textit{audit record} memiliki format yang konsisten. \\
\hline

Audit Event Catalogue  &
Menyediakan definisi, metadata, dan struktur setiap jenis \textit{audit event} yang digunakan oleh sistem. \\
\hline

Audit Repository &
Menyimpan seluruh \textit{audit record} yang dihasilkan sehingga dapat digunakan sebagai bukti pada proses audit maupun investigasi. \\
\hline

Audit Monitor  &
Melakukan pemantauan terhadap aliran \textit{audit event} untuk mendeteksi pola atau kondisi tertentu yang memerlukan perhatian lebih lanjut. \\
\hline

Audit Alert  &
Menghasilkan notifikasi apabila Audit Monitor mendeteksi pola aktivitas yang memenuhi kondisi tertentu sesuai kebijakan yang ditetapkan. \\
\hline

Audit Report &
Menyediakan fasilitas pencarian, pengambilan, dan penyajian \textit{audit record} untuk mendukung proses audit dan investigasi. \\
\hline

\end{xltabular}

Berdasarkan identifikasi tersebut, dapat dilihat bahwa sistem \textit{audit trail} menurut ISO 27789 terdiri atas beberapa layanan yang memiliki tanggung jawab berbeda namun saling melengkapi. Audit Logger Service berperan sebagai titik masuk seluruh \textit{audit event}, sedangkan Audit Repository bertanggung jawab sebagai media penyimpanan utama. Di sisi lain, Audit Report Service menyediakan akses terhadap data audit untuk keperluan investigasi, sementara Audit Monitor Service dan Audit Alert Service mendukung kemampuan pemantauan secara \textit{real-time} terhadap aktivitas yang dicatat. Selain itu, Audit Event Catalogue Service dan Audit Record Generator Service berfungsi menjaga konsistensi struktur \textit{audit record} yang dihasilkan oleh berbagai aplikasi.

Komponen-komponen tersebut selanjutnya akan dibandingkan dengan kondisi eksisting pada arsitektur DecMed untuk mengidentifikasi komponen yang telah tersedia maupun komponen yang masih belum tersedia. Hasil analisis tersebut menjadi dasar dalam menentukan kebutuhan infrastruktur \textit{audit trail} yang akan dirancang pada penelitian ini.